\section{Thảo Luận}
\label{sec:discussion}

Các kết quả được trình bày trong phần trước xác nhận tính hiệu quả, hiệu suất và khả năng thích
nghi của khung DeSFAM trong việc tối thiểu hóa bề mặt tấn công lời gọi hệ thống và phát hiện
hành vi bất thường trong các môi trường container hóa. Phần này thảo luận về những hiểu biết
chính, các đánh đổi thiết kế và các hàm ý thực tế của việc triển khai DeSFAM trong các tình
huống bảo mật container thực tế.

\subsection{Ưu Điểm Của Phân Tích Lai}
\label{subsec:hybrid_advantage}

Một trong những điểm mạnh cốt lõi của DeSFAM nằm ở cách tiếp cận lập hồ sơ lời gọi hệ thống
lai. Bằng cách tích hợp phân tách tĩnh các nhị phân container với theo dõi thời gian chạy động
sử dụng eBPF và bpftrace, DeSFAM nắm bắt cả hành vi lời gọi hệ thống dự kiến lẫn phụ thuộc
vào thời gian chạy. Tầm nhìn hai lớp này đảm bảo rằng các lời gọi hệ thống hiếm nhưng thiết yếu
không bị bỏ sót, không giống như các hệ thống chỉ tĩnh như Confine~\cite{ref31_confine} vốn
gặp vấn đề xấp xỉ dưới, hay các hệ thống chỉ động như NIMOS~\cite{ref38_nimos} có nguy cơ
xấp xỉ trên.

\subsection{Giảm Thiểu Vượt Trội Mà Không Mất Chức Năng}
\label{subsec:reduction_advantage}

DeSFAM đạt được tỷ lệ giảm lời gọi hệ thống trung bình là 78,60\%, cao hơn đáng kể so với
Confine (47,60\%), Prof-gen (67,80\%) và NIMOS (63,40\%). Bất chấp sự giảm thiểu tích cực này,
tất cả các container mục tiêu duy trì chức năng đầy đủ trong suốt tất cả các bài kiểm tra khối
lượng công việc. Điều này xác nhận rằng việc tổng hợp danh sách truy cập của DeSFAM cân bằng
bảo mật và tương thích tốt hơn so với các bộ tạo danh sách truy cập tĩnh hoặc dựa trên nhật ký.

\subsection{Thực Thi Thời Gian Thực Hiệu Năng Cao}
\label{subsec:performance_advantage}

Việc sử dụng eBPF của khung để lọc lời gọi hệ thống và quản lý chính sách giới thiệu chi phí
thời gian chạy tối thiểu. Mức sử dụng CPU tăng trung bình chỉ 1,74\% và mức sử dụng bộ nhớ
tăng khoảng 50~MB, hoàn toàn dưới ngưỡng ảnh hưởng đến hiệu năng container. Độ trễ lời gọi
hệ thống duy trì dưới 300~nano-giây mỗi lời gọi, xác nhận sự phù hợp của DeSFAM cho các khối
lượng công việc thông lượng cao như Nginx và Redis. Để so sánh, các giải pháp không gian người
dùng (ví dụ: auditd) thường giới thiệu độ trễ hàng mili-giây.

\subsection{Phản Ứng Thích Nghi Với Các Mối Đe Dọa}
\label{subsec:adaptive_advantage}

Logic thực thi của DeSFAM không tĩnh. Nó cập nhật động các danh sách cho phép lời gọi hệ thống
bằng cách sử dụng chấm điểm rủi ro và đầu ra phát hiện bất thường từ mô-đun SyscallAD. Mô hình
đạt 96\% độ chính xác tổng thể, với điểm F1 phát hiện tấn công là 0,92 và tỷ lệ dương tính giả
thấp. Điều này cho phép DeSFAM phát hiện và phản ứng với các cuộc tấn công CVE nâng cao (ví dụ:
Polkit và Dirty Pipe) trong thời gian thực, loại bỏ các lời gọi hệ thống dễ bị tổn thương như
\texttt{setuid}, \texttt{splice} và \texttt{clone} ngay khi phát hiện.

\subsection{Hạn Chế Và Hướng Nghiên Cứu Tương Lai}
\label{subsec:limitations}

Mặc dù DeSFAM giảm hiệu quả sự phơi nhiễm lời gọi hệ thống và phát hiện hành vi bất thường, độ
chính xác của nó vẫn nhạy cảm với tính đại diện của dữ liệu huấn luyện. Các nghiên cứu trong
tương lai sẽ tích hợp lập hồ sơ nhận thức khối lượng công việc, phân loại vai trò container và
học liên tục để cải thiện khả năng tổng quát hóa trên các khối lượng công việc đa dạng.

Lọc hiện tại hoạt động ở cấp độ lời gọi hệ thống (dựa trên ID và chuỗi) mà không có kiểm tra
đối số. Bảng~\ref{tab:cve_filter_rules} tóm tắt các quy tắc lọc theo ngữ cảnh được áp dụng
trong quá trình trích xuất lời gọi hệ thống dựa trên CVE. Các mở rộng trong tương lai sẽ khám
phá phân tích cấp độ đối số và theo dõi truy cập bộ nhớ qua eBPF để cải thiện thêm độ chi tiết
phát hiện và khả năng phục hồi trước các cuộc tấn công tinh vi.

\begin{table}[htbp]
  \centering
  \caption{Các quy tắc lọc để trích xuất lời gọi hệ thống dựa trên CVE.}
  \label{tab:cve_filter_rules}
  \includegraphics[width=\linewidth]{table14}
\end{table}

Khả năng mở rộng trong các môi trường điều phối cũng là một trọng tâm. Các triển khai trong
tương lai sẽ tích hợp với Kubernetes qua các hook CRI và thực thi chính sách nhận thức không
gian tên, đảm bảo bảo vệ nhất quán, chi phí thấp trên các cụm container hóa động.
